\documentclass[11pt]{article}
\usepackage[margin=1in]{geometry}
\usepackage[T1]{fontenc}
\usepackage[utf8]{inputenc}
\usepackage{array}
\usepackage{booktabs}
\usepackage{enumitem}
\usepackage{url}
\usepackage{xcolor}
\usepackage{listings}
\def\UrlBreaks{\do\/\do-\do\_\do\.\do\?\do\=\do\&}

\lstdefinestyle{plaincmd}{
  basicstyle=\ttfamily\small,
  columns=fullflexible,
  breaklines=true,
  frame=single,
  backgroundcolor=\color{black!3}
}
\setlength{\emergencystretch}{3em}

\title{ES-FA / BTP Progress Report}
\author{Event-Driven SNN FPGA Accelerator}
\date{Updated \today}

\begin{document}
\maketitle

\section*{Short Summary}
The project has moved beyond the initial proposal stage. At this point, ES-FA is
a local working pipeline with SNN training, hardware-aware estimates, experiment
runs, analysis plots, RTL blocks, xsim/Vivado validation scripts, a small CLI
layer, and LaTeX documentation. The main unfinished item is physical validation
on the KV260 board. Local xsim/Vivado evidence already exists, but board-level
measurements still have to be added.

\section{Project Aim}
The idea behind ES-FA is to study whether sparse spiking neural network activity
can be used to reduce the work done by an FPGA accelerator. The project therefore
tracks both model quality and hardware-facing cost. Accuracy alone is not enough;
the run also records spike activity, memory-access estimates, energy proxy,
latency proxy, and hardware validation data where available.

In simple terms, the work asks:
\begin{itemize}[leftmargin=*]
\item Can a small SNN be trained and exported in a reproducible way?
\item Can dense and event-driven execution be compared using the same metrics?
\item Can the RTL and validation flow produce local evidence for the KV260 path?
\item Can the project be demonstrated as a runnable system instead of a loose set
of scripts?
\end{itemize}

\section{Work Completed So Far}
\begin{center}
\begin{tabular}{p{0.24\linewidth}p{0.64\linewidth}}
\toprule
Area & Status \\
\midrule
Training & Baseline LIF SNN training is available for a \texttt{784 -> 128 -> 64 -> 10} network. \\
Estimators & Spike cost, memory access, energy proxy, and latency proxy functions are connected to training. \\
Experiments & Hardware-aware loss and adaptive dataflow runs are present in the current results. \\
Exports & Metrics, histories, manifests, checkpoints, INT8 weights, spike stats, and raster files are generated. \\
RTL & Router, scheduler, event queue, BRAM, LIF PE, and top-level Verilog sources were preserved. \\
Validation & KV260-oriented xsim and Vivado batch scripts generate local logs, reports, and parsed JSON. \\
System layer & Intent engine, task router, SNN/math/control modules, adaptive policy, and CLI are implemented. \\
Docs & Setup guide, manual, architecture notes, checkpoints, paper source, and this report are now in the repo. \\
\bottomrule
\end{tabular}
\end{center}

\section{Existing Work Preserved}
The project was not rebuilt from scratch. The following parts were kept because
they were already useful and working:
\begin{itemize}[leftmargin=*]
\item \path{paper1_training/}: current SNN model, training core, baseline script,
and export utilities.
\item \path{paper2_hardware_model/estimators.py}: hardware proxy functions.
\item \path{experiments/}: proposal experiment folders and configs.
\item \path{iterations/}: refinement runs.
\item \path{analysis/compare.py}: ranking and plotting workflow.
\item \path{hardware/}: Verilog RTL and testbenches.
\item \path{software/}: earlier train/export/bank-mapping scripts.
\end{itemize}

\section{What Was Added or Improved}
The newer work mainly wraps the existing pieces into a cleaner local pipeline:
\begin{itemize}[leftmargin=*]
\item run manifests and layer-activity exports were added so results are easier
to trace;
\item \path{hardware_validation/kv260/} was added for xsim regression, Vivado Tcl
builds, report parsing, and hardware metric collection;
\item \path{system_layer/} was added for intent classification and routing;
\item \path{deployment/} was added for the CLI and adaptive runtime policy;
\item \path{analysis/evidence_report.py} was added for baseline-vs-optimized
tables and interpretations;
\item presentation-facing LaTeX documents were added under \path{docs/} and
\path{paper_output/}.
\end{itemize}

\section{Current Architecture}
The current repo has four practical layers:
\begin{lstlisting}[style=plaincmd]
Software:
  SNN training -> hardware estimators -> experiments -> analysis

Hardware:
  Verilog RTL -> xsim regression -> Vivado build -> parsed reports

Runtime:
  intent engine -> router -> SNN/math/control module -> adaptive policy

Outputs:
  metrics -> plots -> evidence tables -> LaTeX reports/PDFs
\end{lstlisting}

The SNN model used in the current baseline is:
\begin{lstlisting}[style=plaincmd]
MNIST image: 28 x 28
Flattened input: 784
Hidden layer 1: 128 LIF neurons
Hidden layer 2: 64 LIF neurons
Output layer: 10 classes
Default temporal window: 16 time steps
\end{lstlisting}

The FPGA-side flow being validated is:
\begin{lstlisting}[style=plaincmd]
input events -> spike router -> scheduler/event queue
  -> weight/neuron memory -> LIF processing element
  -> writeback and counters
\end{lstlisting}

\section{Software Results}
The current local ranking compares the baseline with two experiment outputs:
\begin{center}
\begin{tabular}{lrrrr}
\toprule
Run & Accuracy & Sparsity & Energy proxy & Score \\
\midrule
exp1\_hardware\_aware\_loss & 0.9570 & 0.5648 & 4,592,863.17 & 0.7699 \\
baseline\_paper1 & 0.9570 & 0.5648 & 22,604,295.76 & 0.7387 \\
exp5\_dataflow\_adaptation & 0.9570 & 0.5648 & 45,016,894.73 & 0.6699 \\
\bottomrule
\end{tabular}
\end{center}

The strongest current software result is \path{exp1_hardware_aware_loss}. In the
current run, it keeps the same accuracy as the baseline while reducing the energy
proxy by about 79.68 percent against the dense baseline estimate. This is still
a proxy result, so it should not be presented as measured board power.

\section{Hardware Validation}
The hardware validation path is local and KV260-oriented. It runs:
\begin{enumerate}[leftmargin=*]
\item xsim regression for dense and event scheduler modes;
\item Vivado synthesis/implementation through \path{hardware_validation/kv260/vivado/build.tcl};
\item parsing of utilization and timing reports;
\item merging of xsim, Vivado, model, mode, and estimator fields into
\path{hardware_metrics.json}.
\end{enumerate}

Current local hardware evidence exists for baseline and adaptive variants in
dense and event modes. The current parsed results show:
\begin{itemize}[leftmargin=*]
\item latency: 5760 ns in the active validation window;
\item cycle count: 576 cycles;
\item LUT usage: 22,193 in the parsed reports;
\item estimator-vs-measured latency MAPE: 33.33 percent across four hardware runs.
\end{itemize}

The dense and event modes currently report equal latency/cycle count in the
generated evidence table. That result is useful because it keeps the comparison
honest, but it also shows what needs work next: the sparse event path needs
stronger hardware stimuli and real KV260 board measurements.

\section{System and CLI Layer}
The system layer was added so the project can be shown through commands, not only
through training scripts. The CLI currently supports:
\begin{lstlisting}[style=plaincmd]
python deployment/cli_interface/main.py --query "hi"
python deployment/cli_interface/main.py --query "strike rate 167.55 balls 39"
python deployment/cli_interface/main.py --query "classify results/runtime/sample_signal.bin"
\end{lstlisting}

Internally, the intent engine classifies a request as control, math, or SNN
execution. The router sends it to the matching module. For SNN-style requests,
the adaptive policy estimates signal sparsity and chooses dense or event mode.
Every adaptive decision is logged in:
\begin{lstlisting}[style=plaincmd]
results/runtime/adaptive_decisions.jsonl
\end{lstlisting}

In the current sample signal, the adaptive path selects dense mode because the
estimated sparsity is about 0.0039, below the 0.9000 threshold.

\section{Documentation Added}
The documentation now has a more usable shape:
\begin{itemize}[leftmargin=*]
\item \path{docs/first_time_user_setup_guide.tex}: practical setup guide;
\item \path{docs/architecture.md}: current architecture notes;
\item \path{docs/checkpoints.md}: milestone and run-history ledger;
\item \path{docs/manual.tex}: short operating manual;
\item \path{paper_output/progress_report.tex}: this report;
\item \path{paper_output/final_paper.tex}: paper-style draft source.
\end{itemize}

\section{Open Work}
\begin{itemize}[leftmargin=*]
\item Run physical KV260 measurements and merge them into the same hardware
metrics schema.
\item Tune estimator constants using measured hardware data.
\item Expand the hardware validation matrix beyond the current baseline and
\path{exp5_dataflow_adaptation} variants.
\item Improve event-mode test stimuli so sparse activity is exercised more
clearly.
\item Revisit timing closure where Vivado reports negative slack.
\item Keep generated tool outputs out of the repo root when preparing the GitHub
version.
\end{itemize}

\section{Next Steps}
\begin{enumerate}[leftmargin=*]
\item Boot the KV260 with the prepared Kria Ubuntu SD card.
\item Re-run the smoke pipeline before the presentation to refresh visible
outputs.
\item Run xsim-only validation if Vivado implementation time is limited.
\item Run full Vivado validation for baseline and adaptive variants when time is
available.
\item Add board-side measurements to \path{results/hardware_validation/}.
\item Use those measurements to improve the estimator-vs-hardware match.
\end{enumerate}

\section{Final Note for Review}
The project is now in a demonstrable state, but it should be described carefully:
the software pipeline, plots, CLI, RTL, and local validation scripts are working;
the final board measurement stage is still pending. The current evidence supports
the ES-FA direction, especially the hardware-aware software result, but the final
claim should wait for KV260 physical measurements.

\end{document}
